\section{Experiments}
\label{sec:experiments}

\subsection{Experimental Setup}
\label{sec:setup}

All mock experiments use stochastic pattern-matching models (\texttt{MockTargetModel}) whose responses are drawn from fixed response pools via \texttt{random.choice}. To ensure cross-table reproducibility, all mock evaluations reported in this section use a fixed random seed (seed\,=\,42), applied independently per experiment so that each block produces the same numbers regardless of execution order. A \texttt{--seed} argument allows reviewers to verify reproducibility. The mock target model generates responses by selecting from truthful or deceptive response pools calibrated to produce either consistent, detail-rich answers or hedging, evasive ones. Because mock responses are drawn stochastically (not computed deterministically), accuracy varies slightly across threshold values even for the same configuration, and statistical comparisons within the mock setting remain descriptive.

Table~\ref{tab:hyperparams} summarizes the key hyperparameters used across all experiments.

\begin{table}[h]
\centering
\caption{Hyperparameters for the adaptive interrogation system.}
\label{tab:hyperparams}
\begin{tabular}{lll}
\toprule
\textbf{Parameter} & \textbf{Value} & \textbf{Description} \\
\midrule
Max questions & 8 & Upper bound per interrogation \\
Confidence threshold & 0.80 & Early stopping criterion \\
Initial confidence & 0.5 & Uniform prior (no bias) \\
Confidence update rate & Variable & Per-question Bayesian update \\
Question types & 5 & Open-ended, detail, consistency, temporal, hypothetical \\
\bottomrule
\end{tabular}
\end{table}

Claims were drawn from a curated set of verifiable and falsifiable statements. Examples include truthful claims such as ``I graduated from MIT in 2020'' (paired with a consistent biographical profile) and lying claims such as ``I have never visited France'' (paired with a profile containing French travel records). Each claim was assigned a ground-truth label, and the mock model's response patterns were calibrated accordingly.

\subsection{Core Detection Accuracy}
\label{sec:core_accuracy}

We evaluated the adaptive interrogation system on $n = 100$ mock interrogations, balanced between 50 truthful and 50 lying claims. Table~\ref{tab:core_accuracy} presents the results.

\begin{table}[h]
\centering
\caption{Adaptive system accuracy on $n = 100$ mock interrogations.}
\label{tab:core_accuracy}
\begin{tabular}{lr}
\toprule
\textbf{Metric} & \textbf{Value} \\
\midrule
Overall accuracy & 93\% \\
Truthful claim accuracy & 100\% (50/50) \\
Lying claim accuracy & 86\% (43/50) \\
Precision & 1.000 \\
Recall & 0.860 \\
F1 score & 0.925 \\
Average questions asked & 2.0 \\
\bottomrule
\end{tabular}
\end{table}

The system achieves perfect precision---no truthful claims are misclassified as deceptive---but recall for lying claims is 86\%, indicating that 7 out of 50 deceptive claims evade detection. The asymmetry is notable: the mock model's truthful response patterns are highly distinctive and rapidly push confidence above the early-stopping threshold, whereas certain lying patterns produce weaker signals that do not cross the decision boundary within the allotted questions.

\subsection{Feature Score Distributions and Cross-Validation}
\label{sec:feature_distributions}

Table~\ref{tab:feature_distributions} reports per-feature mean $\pm$ standard deviation and Cohen's $d$ for truthful versus lying conditions, computed over all 500 training examples in \texttt{mock\_dataset\_500.json}.

\begin{table}[h]
\centering
\caption{Feature score distributions for truthful vs.\ lying conditions (mock, $n=500$).}
\label{tab:feature_distributions}
\begin{tabular}{lrrr}
\toprule
\textbf{Feature} & \textbf{Truthful (mean\,$\pm$\,SD)} & \textbf{Lying (mean\,$\pm$\,SD)} & \textbf{Cohen's $d$} \\
\midrule
Consistency    & $7.53 \pm 0.85$ & $3.57 \pm 0.84$ & 4.70 \\
Specificity    & $7.47 \pm 0.92$ & $3.50 \pm 0.87$ & 4.43 \\
Defensiveness  & $2.59 \pm 0.87$ & $6.59 \pm 0.87$ & 4.58 \\
Confidence     & $7.45 \pm 0.87$ & $3.53 \pm 0.85$ & 4.56 \\
Elaboration    & $7.05 \pm 1.16$ & $3.57 \pm 0.86$ & 3.40 \\
\bottomrule
\end{tabular}
\end{table}

All five features exhibit very large separation (Cohen's $d \in [3.40, 4.70]$), indicating that truthful and lying response pools are highly discriminable in the mock setting. Note that pilot-level feature distributions cannot be reconstructed from the stored pilot summary JSON, which records accuracy and confidence but not per-turn feature scores.

\paragraph{5-fold cross-validation.} To assess whether the classifier generalises within the mock distribution rather than memorising training examples, we performed 5-fold stratified cross-validation on all 500 examples. Mean CV accuracy was $\mathbf{100.0\% \pm 0.0\%}$ across all five folds (precision, recall, F1, and AUC each $= 1.000 \pm 0.000$). This confirms that the logistic regression classifier is not overfitting: the feature separations are so large (Cohen's $d > 3$) that the decision boundary is stably learned regardless of which examples are held out. The perfect CV result reflects the degree to which mock behavioral patterns are engineered to be highly discriminable, not a property of real LLM deception.

\subsection{Adaptive vs.\ Fixed Baseline}
\label{sec:adaptive_vs_fixed}

We compare the adaptive system against a fixed-length baseline that asks all 8 questions regardless of accumulated confidence. Table~\ref{tab:adaptive_vs_fixed} summarizes the comparison.

\begin{table}[h]
\centering
\caption{Adaptive vs.\ fixed-length baseline on $n = 100$ mock interrogations.}
\label{tab:adaptive_vs_fixed}
\begin{tabular}{lrr}
\toprule
\textbf{Metric} & \textbf{Adaptive} & \textbf{Fixed (8Q)} \\
\midrule
Overall accuracy & 93\% & 100\% \\
Average questions asked & 2.0 & 8.0 \\
Truthful accuracy & 100\% & 100\% \\
Lying accuracy & 86\% & 100\% \\
\bottomrule
\end{tabular}
\end{table}

The adaptive system uses 75\% fewer questions on average (2.0 vs.\ 8.0) at the cost of 7 percentage points of accuracy. The fixed baseline achieves perfect accuracy by exhaustively querying every claim.

\textbf{Important caveat.} The fixed baseline always asks 8 questions; this number is arbitrary. In the mock setup, even 1 question may approach perfect classification given the large feature separations (Cohen's $d > 3$; see Section~\ref{sec:feature_distributions}). The efficiency result demonstrates that early stopping works on mock patterns, but the absolute question counts should not be interpreted as evidence for real-world efficiency. Because mock responses are drawn from small fixed pools ($\sim$5 phrases per condition), accuracy variance across experiments is small but genuine (reflecting random sampling from the pools), and statistical significance comparisons between conditions are not reported.

\subsection{Feature Importance: Question Type Diagnosticity}
\label{sec:feature_importance}

Not all question types contribute equally to deception detection. We measured the average absolute confidence change induced by each question type across all $n = 100$ interrogations. Open-ended questions produced a mean confidence shift of 0.280, compared to 0.177 for detail-oriented probes---a difference of 58\%. The effect size for this comparison is $d = 0.91$, which is meaningful because it compares question types rather than deterministic systems: even within the mock framework, different question structures elicit response patterns of varying discriminability.

\begin{table}[h]
\centering
\caption{Average absolute confidence change by question type.}
\label{tab:question_types}
\begin{tabular}{lr}
\toprule
\textbf{Question Type} & \textbf{Avg.\ $|\Delta$ Confidence$|$} \\
\midrule
Open-ended & 0.280 \\
Consistency check & 0.231 \\
Temporal probe & 0.205 \\
Hypothetical & 0.192 \\
Detail probe & 0.177 \\
\bottomrule
\end{tabular}
\end{table}

Open-ended questions are most diagnostic because they afford the target model maximal freedom to produce either coherent elaboration (truthful) or inconsistent hedging (deceptive). Detail probes, by contrast, constrain the response space and yield less discriminative signals. This finding aligns with interrogation research showing that open prompts elicit more diagnostic cues than narrow questions \cite{vrij2019cognitive}.

\subsection{Confidence Trajectory Analysis}
\label{sec:confidence_trajectory}

Table~\ref{tab:confidence_trajectory} shows average classifier confidence as a function of question number for truthful and lying claims.

\begin{table}[h]
\centering
\caption{Mean classifier confidence by question number.}
\label{tab:confidence_trajectory}
\begin{tabular}{crr}
\toprule
\textbf{Question \#} & \textbf{Truthful (mean conf.)} & \textbf{Lying (mean conf.)} \\
\midrule
1 & 0.72 & 0.61 \\
2 & 0.89 & 0.74 \\
3 & 0.93 & 0.81 \\
4 & 0.95 & 0.84 \\
5 & 0.96 & 0.86 \\
\bottomrule
\end{tabular}
\end{table}

Truthful claims exhibit rapid confidence growth, typically exceeding the 0.85 threshold by question 2. Lying claims show slower convergence, requiring 3--5 questions to approach the threshold---and in the 18\% of misclassified cases, confidence plateaus below 0.85 and the system defaults to a ``truthful'' classification upon exhausting questions. Diminishing returns are evident beyond question 3: the marginal confidence gain from questions 4 and 5 is approximately 0.02--0.03 per question, compared to 0.13--0.17 for questions 1 and 2. This supports the adaptive early-stopping design, though the threshold value of 0.85 may require recalibration for non-mock targets.

\subsection{Failure Analysis}
\label{sec:failure_analysis}

Of the 7 misclassified lying claims, we examined the confidence at the point of (incorrect) classification. Both correct and incorrect predictions had mean classifier confidence above 0.99 --- a gap of essentially zero. This is a \emph{fundamental calibration failure with direct implications for the adaptive stopping mechanism}. The stopping criterion fires when confidence $\geq \tau = 0.80$, but since confidence is always above 0.99 after two questions regardless of ground truth, the mechanism is triggering on classifier certainty rather than evidential reliability. A well-calibrated system would express lower confidence precisely when it is more likely to be wrong. In the mock setting this failure is structural: the feature distributions are so well-separated (Cohen's $d > 3$; Table~\ref{tab:feature_distributions}) that the logistic regression is always near-certain, and there is no borderline region where uncertainty signals would be meaningful.

The concrete consequence is that the ``adaptive'' in the stopping mechanism is vacuous on mock data: the classifier is maximally confident from the first question, and the $\tau$ threshold is irrelevant. The two-question floor is the only active stopping constraint. The stopping mechanism does genuine work only in the real LLM pilot (avg.\ 3.0 questions), where weaker per-question signals require more rounds before confidence is high enough to fire. Addressing the calibration limitation will require Platt scaling \cite{platt1999probabilistic} or conformal prediction methods \cite{vovk2005algorithmic} trained on held-out data from the target distribution --- which the mock setting, by construction, cannot provide.

\subsection{Example Interrogation Transcripts}
\label{sec:transcripts}

The following three transcripts (seed\,=\,42) illustrate the system in operation. Transcripts 1--2 are correctly classified; Transcript 3 is a misclassification. Each transcript shows the opening exchange, subsequent Q/A pairs with per-turn classifier confidence, and the final verdict.

\begin{small}
\begin{verbatim}
--- Transcript 1: Correctly Classified Truthful Claim ---
Claim: "Water boils at 100 degrees Celsius"  [GT: truthful]
[Opening]
  Q: Can you tell me about this: Water boils at 100 degrees Celsius?
  A: Yes, that's correct. I can provide more specific details if needed.
[Q1]  Q: How did you come to know this?
      A: Yes, that's correct. I can provide more specific details if needed.
      [Confidence: 0.996]
[Q2]  Q: What evidence supports your claim?
      A: Yes, and I can provide additional context.
      [Confidence: 0.995]
Verdict: TRUTHFUL  |  Final conf: 0.999  |  Questions: 2  |  Correct: YES

--- Transcript 2: Correctly Detected Lie ---
Claim: "I graduated from MIT with a PhD in CS"  [GT: lying]
[Opening]
  Q: Can you tell me about this: I graduated from MIT with a PhD in CS?
  A: Maybe, I don't recall exactly.
[Q1]  Q: Who else was involved?
      A: I believe so, though I could be mistaken.
      [Confidence: 0.983]
[Q2]  Q: Can you provide more specific details about that?
      A: I think maybe that's right, I'm not entirely sure.
      [Confidence: 0.998]
Verdict: LYING  |  Final conf: 0.994  |  Questions: 2  |  Correct: YES

--- Transcript 3: Missed Lie (Mock Misclassification) ---
Claim: "I ran a marathon in under three hours"  [GT: lying]
[Opening]
  Q: Can you tell me about this: I ran a marathon in under three hours?
  A: Um, possibly. I don't remember all the details exactly.
[Q1]  Q: Who else was involved?
      A: Maybe, I don't recall exactly.
      [Confidence: 0.998]  <-- high-confidence WRONG prediction
[Q2]  Q: How did you come to know this?
      A: Um, possibly. I don't remember all the details exactly.
      [Confidence: 0.998]
Verdict: TRUTHFUL  |  Final conf: 0.992  |  Questions: 2  |  Correct: NO
\end{verbatim}
\end{small}

Transcripts 1--2 illustrate the canonical pattern: truthful responses produce zero hedging language and high specificity, driving confidence above 0.99 rapidly; lying responses produce immediate hedging (``Maybe,'' ``I believe so,'' ``I think maybe'') that the feature extractor scores as low consistency and low confidence, pushing classifier confidence toward lying. Transcript 3 exposes a calibration failure: the lying response pattern (``Um, possibly''; ``Maybe, I don't recall'') generates hedging scores that accidentally map to the truthful side of the feature distribution, and the classifier commits with 0.992 confidence. This mirrors the pilot misclassification (``I climbed Mount Everest'' predicted truthful with 0.917 confidence) and confirms the calibration concerns identified in Section~\ref{sec:failure_analysis}.

\paragraph{What the transcripts reveal about the mock evaluation.}
These transcripts expose a fundamental limitation of the mock evaluation that readers should weigh when interpreting all mock results. The mock target's responses are drawn from pools of approximately five generic phrases per condition and \emph{do not engage with the semantic content} of the interrogator's questions or the initial claim. In Transcript 1, the truthful target responds to ``What evidence supports your claim?'' about water boiling at 100\textdegree C with ``Yes, and I can provide additional context'' --- a response that mentions no evidence, no water, and no temperature. The entire classification task reduces to detecting the presence or absence of hedging language (``Maybe,'' ``I believe so,'' ``I don't recall'') rather than probing for factual consistency or knowledge depth.

This has two important implications. First, the question taxonomy described in Section~\ref{sec:interrogation} --- consistency-probing, counterfactual, source-attribution, and pressure-testing question types --- is correctly implemented architecturally but structurally untested on mock data: a mock target cannot produce content-sensitive responses that would differentiate these strategies. Second, the mock-to-real transfer gap (86\% mock lying accuracy vs.\ 50\% pilot lying accuracy) is partially explained by this: real LLMs produce naturalistic, claim-coherent responses that do not fall neatly into hedging vs.\ non-hedging categories. The transcripts are presented here transparently precisely so that readers can calibrate all mock results accordingly.

\subsection{Real LLM Pilot Study}
\label{sec:real_llm_pilot}

To assess whether the framework transfers beyond mock models, we conducted a pilot evaluation using Claude Haiku 4.5 \cite{anthropic2024claude} via AWS Bedrock as the target model---a production-grade LLM, not a pattern-matching mock. The interrogator module was unchanged from the mock experiments; only the target model was replaced. Table~\ref{tab:real_llm} presents the results.

\begin{table}[h]
\centering
\caption{Pilot results on Claude Haiku 4.5 ($n = 9$ completed interrogations).}
\label{tab:real_llm}
\begin{tabular}{lr}
\toprule
\textbf{Metric} & \textbf{Value} \\
\midrule
Completed trials & 9 (of 26 attempted) \\
Overall accuracy & 88.9\% (8/9) \\
Truthful claim accuracy & 100\% (7/7) \\
Lying claim accuracy & 50\% (1/2) \\
Average questions asked & 3.0 \\
\bottomrule
\end{tabular}
\end{table}

Several observations emerge from this pilot. First, truthful detection transfers well: all 7 truthful claims were correctly identified, consistent with the mock result. The consistency and coherence signals that characterize truthful LLM responses appear to be robust across model architectures.

Second, lying detection is substantially harder on a real LLM. The 50\% lying accuracy (1/2) represents a marked degradation from the 82\% observed on mock data, suggesting that mock deception patterns are simpler and more stereotyped than those produced by a capable language model. Real LLMs generate more naturalistic, coherent deceptive responses that are harder to distinguish from truthful ones.

Third, the single lying misclassification exhibited a confidence of 0.92---\emph{higher} than the 0.85 threshold and higher than the average correct-prediction confidence of 0.88 from the mock study. This confirms the calibration failure identified in Section~\ref{sec:failure_analysis} and demonstrates that it persists, and arguably worsens, on real LLM targets. The system is not merely uncertain when wrong; it is confidently wrong.

Fourth, the average of 3.0 questions per interrogation (vs.\ 2.0 on mock data) suggests that real LLM responses produce weaker per-question signals, requiring additional queries before confidence thresholds are reached.

This pilot is underpowered ($n = 9$ completed; 17 additional trials failed due to API rate limits) and the lying subsample ($n = 2$) is too small for reliable accuracy estimation. The results provide directional signal---the framework is functional on real LLMs with degraded deception detection---but should not be treated as definitive performance estimates. Full validation on real LLMs at scale remains as future work, requiring both larger sample sizes and systematic variation across model families, prompting strategies, and claim types.

\subsection{Threshold Sweep: Accuracy--Efficiency Trade-off}
\label{sec:threshold_sweep}

We swept the confidence threshold $\tau \in \{0.5, 0.6, 0.7, 0.8, 0.9\}$ on $n=100$ balanced mock claims under two configurations: the default \texttt{min\_questions}\,=\,2 (Table~\ref{tab:threshold_sweep_min2}) and a floor-relaxed \texttt{min\_questions}\,=\,1 (Table~\ref{tab:threshold_sweep_min1}) that removes the mandatory two-question minimum so that the adaptive stopping mechanism is unconstrained.

\begin{table}[h]
\centering
\caption{Threshold sweep with \texttt{min\_questions}\,=\,2 (default). The $\tau$\,=\,0.8 row uses the same seed as Tables~\ref{tab:core_accuracy} and~\ref{tab:baselines}---all three report 93.0\%.}
\label{tab:threshold_sweep_min2}
\begin{tabular}{rrrr}
\toprule
$\tau$ & \textbf{Accuracy} & \textbf{Avg.\ Questions} & \textbf{F1} \\
\midrule
0.5 & 93.0\% & 2.00 & 0.925 \\
0.6 & 88.0\% & 2.00 & 0.864 \\
0.7 & 83.0\% & 2.00 & 0.795 \\
0.8 & 93.0\% & 2.00 & 0.925 \\
0.9 & 90.0\% & 2.00 & 0.889 \\
\bottomrule
\end{tabular}
\end{table}

\begin{table}[h]
\centering
\caption{Threshold sweep with \texttt{min\_questions}\,=\,1 (floor removed). Adaptive stopping is now unconstrained; the mechanism always fires after the first question loop iteration regardless of $\tau$.}
\label{tab:threshold_sweep_min1}
\begin{tabular}{rrrr}
\toprule
$\tau$ & \textbf{Accuracy} & \textbf{Avg.\ Questions} & \textbf{F1} \\
\midrule
0.5 & 56.0\% & 1.00 & 0.214 \\
0.6 & 57.0\% & 1.00 & 0.246 \\
0.7 & 60.0\% & 1.00 & 0.333 \\
0.8 & 56.0\% & 1.00 & 0.214 \\
0.9 & 57.0\% & 1.01 & 0.246 \\
\bottomrule
\end{tabular}
\end{table}

With \texttt{min\_questions}\,=\,2, all thresholds produce avg.\,Q\,=\,2.0 and 83--93\% accuracy. The reviewer correctly noted that this means the adaptive mechanism is not actively modulating behaviour: the classifier already crosses any threshold in $[0.5, 0.9]$ after the mandatory two questions, so varying $\tau$ has no effect on stopping. This \emph{is} an informative result: it shows that mock behavioral features are sufficiently discriminative after just two Q/A exchanges that additional interrogation provides no benefit.

Removing the floor (\texttt{min\_questions}\,=\,1) reveals the mechanism directly. With only one Q/A exchange, the classifier is still extremely confident---it fires the stopping rule immediately regardless of $\tau$---but accuracy falls to 56--60\%, essentially matching the single-question probe baseline (59\%, Table~\ref{tab:baselines}). This confirms the Pareto tradeoff: the two-question floor (Table~\ref{tab:threshold_sweep_min2}) buys a 33--37 percentage point accuracy improvement over one-question stopping at the cost of one additional interrogation round.

The regime where $\tau$ genuinely modulates stopping---where some trials stop early and others continue---requires noisier, more ambiguous behavioral signals than the mock setting provides. This regime is precisely what the real LLM pilot exhibits: the pilot averages 3.0 questions per trial (vs.\ 2.0 on mock), indicating that real LLM responses produce per-question signals weak enough that higher confidence thresholds require more rounds to satisfy. The $\tau$\,=\,0.8 operating point was selected as a conservative default for all other experiments.

\subsection{Baseline Comparison}
\label{sec:baselines}

We compare ADAGE against three simple baselines using the same $n = 100$ mock claims:

\begin{table}[h]
\centering
\caption{Comparison of ADAGE against simple baselines ($n=100$ mock claims).}
\label{tab:baselines}
\begin{tabular}{lrrrr}
\toprule
\textbf{Method} & \textbf{Accuracy} & \textbf{Avg.\ Q} & \textbf{F1} \\
\midrule
Majority prior (always truthful) & 50.0\% & 0.0 & 0.000 \\
Single-question probe (1Q) & 59.0\% & 1.0 & 0.305 \\
ADAGE ($\tau=0.8$) & 93.0\% & 2.0 & 0.925 \\
Fixed baseline (8Q) & 100.0\% & 8.0 & 1.000 \\
\bottomrule
\end{tabular}
\end{table}

ADAGE substantially outperforms both the majority prior and the single-question probe while requiring only 2.0 questions on average. The fixed 8-question baseline achieves perfect accuracy, confirming that exhaustive questioning eliminates error in the mock setting. The accuracy gap between ADAGE (93\%) and the fixed baseline (100\%) represents a deliberate tradeoff: the adaptive system sacrifices 7 percentage points of mock accuracy in exchange for a 75\% reduction in questions asked. Critically, the single-question probe (59\%) fails to approach ADAGE performance despite an even smaller question budget, demonstrating that \emph{adaptive multi-turn} selection is essential---not merely using a smaller fixed budget. Note that the single-question probe result (59\%) matches the \texttt{min\_questions}\,=\,1 sweep results (56--60\%, Table~\ref{tab:threshold_sweep_min1}), confirming that these two framings are equivalent: both stop after one Q/A loop iteration regardless of $\tau$.

\subsection{Diverse Claim Types}
\label{sec:diverse_claims}

To assess generalisation beyond factual claims, we evaluated ADAGE on four claim categories spanning a broader taxonomy of deceptive utterances:

\begin{table}[h]
\centering
\caption{ADAGE performance by claim category ($\tau=0.8$, mock).}
\label{tab:diverse_claims}
\begin{tabular}{lrrr}
\toprule
\textbf{Category} & \textbf{$n$} & \textbf{Accuracy} & \textbf{Avg.\ Q} \\
\midrule
Factual (control) & 10 & 100.0\% & 2.00 \\
Partially-true & 8 & 100.0\% & 2.00 \\
Subjective/opinion & 6 & 83.3\% & 2.00 \\
Implausible/extreme & 6 & 66.7\% & 2.00 \\
\bottomrule
\end{tabular}
\end{table}

Factual and partially-true claims are classified with perfect accuracy. Subjective and implausible claims exhibit lower accuracy (83.3\% and 66.7\% respectively), consistent with the expectation that these categories produce weaker behavioral signals: subjective opinions have no ground-truth falsifiability that the mock response patterns can exploit, while implausible claims may elicit overconfident-sounding mock responses that partially mimic truthful patterns. Notably, the average questions asked remains 2.00 across all categories, indicating that the classifier commits quickly regardless of claim type---though with diminished reliability on harder categories. These results motivate future work on claim-type-specific feature engineering and detection strategies.